\documentclass[11pt,letterpaper]{article}
\usepackage[utf8]{inputenc} %Codificacion del texto (ISO Latin1 encoding)

\usepackage{fancyhdr} %Permite acomodar a tu gusto la parte de arriba y
\usepackage[spanish]{babel} %Permite definir el idioma del dcumento
\usepackage{graphicx} %Permite exportar imagenes en formato eps
\usepackage{url} %Tipo de fuente para correos y paginas
\usepackage{hyperref}
\usepackage{pgf}
\usepackage{fleqn}
\usepackage{amssymb}
\usepackage{fancyvrb}
\usepackage{sectsty}
\usepackage{makeidx}
\usepackage{colortbl} %Permite colocar colores a las tablas
\usepackage{booktabs}
%Margenes%
\parskip 1mm %Espacio entre parrafos

\setlength{\topmargin}{0pt}

\oddsidemargin	0.5cm  % Ancho Letter 21,59cm
\evensidemargin 0.5cm  % Alto  Letter 27,81cm
\textwidth	15.5cm
\textheight	21.0cm
\headsep	4 mm
\parindent	0.5cm
\pagestyle{fancyplain}
\lhead{Taller de Gestión de Proyectos Informáticos} %Parte superior izquierda
\rhead{\bf \it Sintesis} %Parte superior derecha
\lfoot{\it Principios y Conceptos de la Gestión de Proyectos }
\cfoot{} %Parte inferior central
\rfoot{\bf \thepage} %Parte inferior derecha
\renewcommand{\footrulewidth}{0.4pt} %Linea de separacion inferior

\newtheorem{theorem}{Theorem}
\newtheorem{acknowledgement}[theorem]{Acknowledgement}
\newtheorem{algorithm}[theorem]{Algorithm}
\newtheorem{axiom}[theorem]{Axiom}
\newtheorem{case}[theorem]{Case}
\newtheorem{claim}[theorem]{Claim}
\newtheorem{conclusion}[theorem]{Conclusion}
\newtheorem{condition}[theorem]{Condition}
\newtheorem{conjecture}[theorem]{Conjecture}
\newtheorem{corollary}[theorem]{Corollary}
\newtheorem{criterion}[theorem]{Criterion}
\newtheorem{definition}[theorem]{Definition}
\newtheorem{example}[theorem]{Example}
\newtheorem{exercise}[theorem]{Exercise}
\newtheorem{lemma}[theorem]{Lemma}
\newtheorem{notation}[theorem]{Notation}
\newtheorem{problem}[theorem]{Problem}
\newtheorem{proposition}[theorem]{Proposition}
\newtheorem{remark}[theorem]{Remark}
\newtheorem{solution}[theorem]{Solution}
\newtheorem{summary}[theorem]{Summary}
\newenvironment{proof}[1][Proof]{\noindent\textbf{#1.} }{\ \rule{0.5em}{0.5em}}

\newcommand{\primaria}[1]{
	\textbf{\underline{#1}}
}

\newcommand{\foranea}[1]{
	\textbf{\textsl{#1}}
}

\newcommand{\primyfor}[1]{
	\underline{\foranea{#1}}
}

\makeatletter
\newcommand\subsubsubsection{\@startsection {paragraph}{1}{\z@}%
                                   {-3.5ex \@plus -1ex \@minus -.2ex}%
                                   {1.5ex \@plus.2ex}%
                                   {\normalfont\bfseries}}
\newcommand\subsubsubsubsection{\@startsection {subparagraph}{1}{\z@}%
                                   {-3.5ex \@plus -1ex \@minus -.2ex}%
                                   {1.5ex \@plus.2ex}%
                                   {\normalfont\bfseries}}


\makeatother
%\makeindex
\bibliographystyle{plain}

\begin{document}

%%%%%%%%%%%%%%%%%%%%%%%%%%
%Definicion de la portada%
%%%%%%%%%%%%%%%%%%%%%%%%%%
\begin{titlepage}
    \begin{center}
	\begin{tabular}{ccc}
	    \includegraphics[width=3cm]{img/utfsm}
	    & 
	    \hspace{-0.2cm}
	    \begin{tabular}{c}
		Universidad Técnica Federico Santa María \\ \hline
		\vspace{0.2cm}
		Departamento de Informática\\
		\vspace{1.2cm}
	    \end{tabular}
	    \hspace{0.2cm}
	    &
            \includegraphics[width=2cm]{img/di}
	\end{tabular}

	\vspace{1cm}
	%Titulo del Documento
	\begin{tabular}{c}
		\huge{\sc{Entrevista Apreciativa - Parte 2}}\\\\
		\huge{\sc{Empresa Elian y Cía., Ltda.}}\\\\
		%\includegraphics[scale=0.7]{img/portada} \\\\
	\end{tabular}

    \vspace{3cm}
	\begin{tabular}{c}
		\large{\sc{Integrantes}}
	\end{tabular}
	\\
	\begin{tabular}{cr}
         	\normalsize{Rodrigo Fernández Gaete}  & \url{rfernand@csrg.inf.utfsm.cl}\\
         	\normalsize{Cristián Maureira Fredes} & \url{cmaureir@csrg.inf.utfsm.cl}\\
         	\normalsize{Gabriel Zamora Nelson} & \url{gzamora@csrg.inf.utfsm.cl}\\
	\end{tabular}

\vspace{5cm}

	\begin{tabular}{ll}
         	\normalsize \textbf{Asignatura:} & {Comportamiento Organizacional}\\
            \normalsize \textbf{Profesor:}   & {Lautaro Guerra}\\
         	\normalsize \textbf{Fecha de Entrega:} & {14 de Septiembre del 2010}\\
	\end{tabular}
	%Fecha
    %\normalsize{\sc{\today}}\\
    \end{center}
\end{titlepage}

\tableofcontents
\newpage

%\section{Introducción}
%\label{sec:introduccion}
%Una expresión regular,
a menudo llamada también patrón,
es una expresión que describe un conjunto de cadenas sin enumerar sus elementos.

Además,
normalmente representan otro grupo de caracteres mayor,
de tal forma que podemos comparar el patrón con otro conjunto de caracteres para ver las coincidencias.

La mayoría de las formalizaciones proporcionan los siguientes constructores:
una expresión regular es una forma de representar a los lenguajes regulares
(finitos o infinitos) y se construye utilizando caracteres del alfabeto
sobre el cual se define el lenguaje.

Específicamente,
las expresiones regulares se construyen utilizando los operadores
unión, concatenación y clausura de Kleene.

Las expresiones regulares en UNIX (RE's) están recogidas en el estándar POSIX 1003.2.



\section{Administración del plan de riesgo}

\begin{frame}
\frametitle{Administración del plan de riesgo}
Es el proceso de definir como realizar las actividades de gestión de riesgo del proyecto.

Se debe:
\begin{itemize}
    \item Planificación \textbf{cuidadosa} y bien definida mejora las probabilidades de
    triunfar en los 5 otros procesos de la administración del riesgo.
    \item Proceso que \textbf{comienza} cuando el proyecto es concebido.
\end{itemize}
\begin{columns}
	\begin{column}{0.6\textwidth}
	\end{column}
	\begin{column}{0.4\textwidth}
		\includegraphics[width=3cm]{img/random_img_1}
	\end{column}
\end{columns}
\end{frame}

%\subsection{Valores de entrada}
\begin{frame}
\frametitle{Administración del plan de riesgo}
\framesubtitle{Valores de entrada}

\begin{columns}
	\begin{column}{0.8\textwidth}
\begin{itemize}
    \item<1-> \textbf{Alcances} del proyecto.
    \item<2-> Plan de la administración de \textbf{costos}.
    \item<3-> Plan de administración de \textbf{tiempos}.
    \item<4-> Plan de administración de \textbf{Comunicaciones}.
    \item<5-> \textbf{Factores} del ambiente empresarial.
    \item<6-> Información organizacional de los \textbf{procesos}.
\end{itemize}
	\end{column}
	\begin{column}{0.2\textwidth}
		\includegraphics[width=2cm]{img/input}
	\end{column}
\end{columns}
\end{frame}

%\subsection{Técnicas y Herramientas}
\begin{frame}
\frametitle{Administración del plan de riesgo}
\framesubtitle{Técnicas y Herramientas}
\begin{columns}
	\begin{column}{0.8\textwidth}
\begin{itemize}
    \item \textbf{Reuniones} de análisis y elaboración del plan de administración de riesgo.
    \begin{itemize}
        \item<1-> Requieren de todos los \textbf{involucrados} del proyecto.
        \item<2-> Se busca definir las \textbf{actividades} para la administración del riesgo.
        \item<3-> Se calculan los \textbf{costos} respectivos.
        \item<4-> Revisiones de las \textbf{medidas} de contingencia.
        \item<5-> Asignación de \textbf{responsabilidades}.
        \item<6-> \textbf{Categorización} de riesgos por nivel, probabilidad, impacto, etc.
        \item<7-> Generación de \textbf{templates}.
    \end{itemize}
\end{itemize}
	\end{column}
	\begin{column}{0.2\textwidth}
		\includegraphics[width=2cm]{img/tools}
	\end{column}
\end{columns}
\end{frame}

%\subsection{Valores de Salida}
\begin{frame}
\frametitle{Administración del plan de riesgo}
\framesubtitle{Valores de Salida}
\begin{columns}
	\begin{column}{0.8\textwidth}
\begin{itemize}
    \item<1-> \textbf{Metodología}.
    \item<2-> \textbf{Roles} y responsabilidades.
    \item<3-> \textbf{Presupuesto}.
    \item<4-> Periodos de \textbf{evaluación}.
    \item<5-> \textbf{Categorías} de riesgos.
    \item<6-> \textbf{Definiciones} de probabilidad y impacto de los riesgos.
    \item<7-> Matriz de \textbf{probabilidad} e \textbf{Impacto}.
    \item<8-> \textbf{Tolerancia} de los stakeholders.
    \item<9-> \textbf{Formatos} de reporte.
    \item<10-> Documentación y \textbf{seguimiento}.
\end{itemize}
	\end{column}
	\begin{column}{0.2\textwidth}
		\includegraphics[width=2cm]{img/output}
	\end{column}
\end{columns}
\end{frame}

\section{Identificar Riesgos}
Es el proceso de identificar los riesgo que puedan afectar al proyecto,
documentando sus características principales. Esto es iterativo, ya que nuevos
riesgos pueden evolucionar o aparecer en el transcurso del proyecto. Por ello,
debe involucrar a los miembros del proyecto para inculcar un sentido de
responsabilidad en ellos.

En este proceso, se deben considerar los planes realizados anteriormente,
junto con los costos y duración de cada una de las actividades propuestas y el
lineamiento base sobre los alcances del proyecto. Por ello, los registros con
información sobre los stakeholders y los planes de administración de costos,
tiempos y calidad, junto con los documentos del proyecto, deben ser revisados
exhaustivamente para poder identificara los principales riesgos.

En esta etapa, se debería revisar la documentación que se tenga del proyecto,
y utilizar técnicas de recopilación de la información como Brainstorming,
técnica de Delphi, entrevistas, etc. para poder contar con una buena
información de todos los niveles de trabajo del proyecto. El realizar análisis
con listas de chequeo, analizar las suposiciones y el desarrollo de diagramas
de causa y efecto, flujo del sistema y influencia pueden ayudar a la
identificación de riesgos. Por último, vale la pena nombrar el análisis SWOT
(FODA) y el juicio de expertos para una revisión más profunda de los riesgos.

El resultado de la identificación de riesgos es una documentación con el
registro de riesgos, identificados apropiadamente y con acciones potenciales
para la mitigación de cada uno de ellos.
\subsection{Valores de Entrada}

\begin{enumerate}
    \item Plan de administración de Riesgos.
    \item Costos estimados de las actividades.
    \item Duración estimada de las actividades.
    \item Lineamiento base sobre los alcances del proyecto.
    \item Registros de los stakeholders.
    \item Plan de administración de Costos.
    \item Plan de administración de Tiempos.
    \item Plan de administración de Calidad.
    \item Documentos del proyecto
    \begin{itemize}
        \item Diagramas de red, lineamientos base, eficiencia de trabajo, etc.
    \end{itemize}
    \item Factores del ambiente empresarial.
    \begin{itemize}
        \item Información pública, bases de datos, estudios académicos,
        benchmarking, estudios de la industria, etc.
   
    \end{itemize}
    \item Información organizacional de los procesos.
    \begin{itemize}
        \item Información del proyecto, procesos de control del proyecto,
        templates, lecciones aprendidas, etc.
    \end{itemize}

\end{enumerate}


\subsection{Técnicas y Herramientas}

\begin{enumerate}
    \item Revisiones de la documentación.
    \item Técnicas para la recopilación de la información.
    \begin{itemize}
        \item Brainstorming.
        \item Técnica de Delphi.
        \item Entrevistas.
        \item Análisis de las fuentes del problema.
    \end{itemize}
    \item Análisis con listas de chequeo.
    \item Análisis de las suposiciones.
    \item Técnicas de desarrollo de diagramas.
    \begin{itemize}
        \item Diagramas de causa y efecto, flujo del sistema y influencia.
    \end{itemize}
    \item Análisis SWOT (FODA).
    \item Juicio de Expertos.
\end{enumerate}

\subsection{Valores de Salida}
\begin{enumerate}
    \item Registro de riesgos
    \begin{itemize}
        \item Lista de riesgos identificados.
        \item Lista de acciones potenciales.
    \end{itemize}
\end{enumerate}

\section{Análisis cualitativas del riesgo}
\frame
{
\frametitle{Análisis Cualitativo de Riesgos}
\begin{itemize}
	\item Evalúa la \textbf{prioridad} de los riesgos identificados usando:
	\begin{itemize}
		\item Probabilidad de ocurrencia.
		\item El \textbf{impacto} correspondiente sobre los objetivos del
			proyecto si los riesgos efectivamente ocurren, tales como:
		\begin{itemize}
			\item Costos
			\item Cronograma
			\item Alcance
			\item Calidad
		\end{itemize}
	\end{itemize}
\end{itemize}
\begin{columns}
	\begin{column}{0.6\textwidth}
	\end{column}
	\begin{column}{0.4\textwidth}
		\includegraphics[width=3cm]{img/quality}
	\end{column}
\end{columns}
}

\frame
{
\frametitle{Análisis Cualitativo de Riesgos}
\begin{block}{Definición}
	El Análisis Cualitativo de Riesgos es normalmente una forma rápida de
	\textbf{establecer prioridades} para la Planificación de la Respuesta a los Riesgos.
\end{block}
}


\frame
{
\frametitle{Análisis Cualitativo de Riesgos}
\framesubtitle{Entradas}
\begin{columns}
	\begin{column}{0.8\textwidth}
		\begin{itemize}
			\item<1-> Registro de Riesgos.
			\item<2-> \textbf{Plan} de administración de Riesgos.
			\item<3-> Declaración del \textbf{alcance} del Proyecto.
			\item<4-> Activos organizacionales del Proceso.
		\end{itemize}
	\end{column}
	\begin{column}{0.2\textwidth}
		\includegraphics[width=2cm]{img/input}
	\end{column}
\end{columns}
}

\frame
{
\frametitle{Análisis Cualitativo de Riesgos}
\framesubtitle{Herramientas y Técnicas}
\begin{columns}
	\begin{column}{0.8\textwidth}
		\begin{itemize}
			\item<1-> Probabilidad de riesgos y \textbf{evaluación} de impactos.
			\item<2-> \textbf{Matriz} de probabilidad e impacto.
			\item<3-> Evaluación de la calidad de los datos de los riesgos.
			\item<4-> \textbf{Categorización} de Riesgos.
			\item<5-> Evaluación de la \textbf{urgencia} de los riesgos.
			\item<6-> Juicio experto.
		\end{itemize}
	\end{column}
	\begin{column}{0.2\textwidth}
		\includegraphics[width=2cm]{img/tools}
	\end{column}
\end{columns}
}

\frame
{
\frametitle{Análisis Cualitativo de Riesgos}
\framesubtitle{Salidas}
\begin{columns}
	\begin{column}{0.8\textwidth}
		Actualización de los registros de riesgos.
		\begin{itemize}
			\item<1-> Lista de prioridades de los riesgos del proyecto.
			\item<2-> Riesgos agrupados por \textbf{categorías}.
			\item<3-> Causas de riesgos o \textbf{áreas} que requieren una atención particular.
			\item<4-> Lista de riesgos que requieren respuesta \textbf{cercana} al equipo de trabajo.
			\item<5-> Lista de riesgos para un análisis y respuesta adicional.
			\item<6-> Lista de seguimiento de riesgos de baja prioridad.
			\item<7-> \textbf{Tendencias} en el resultado del análisis de riesgos cualitativo.
		\end{itemize}
	\end{column}
	\begin{column}{0.2\textwidth}
		\includegraphics[width=2cm]{img/output}
	\end{column}
\end{columns}
}

\section{Análisis cuantitativos del riesgo}
\frame
{
\frametitle{Análisis Cuantitativo de Riesgos}
\begin{itemize}
	\item Este proceso usa \emph{técnicas y métodos} \textbf{estadísticos} para:
	\begin{itemize}
		\item Cuantificar los \emph{posibles resultados} del proyecto y sus probabilidades
		\item Evaluar la \emph{probabilidad de lograr} los \textbf{objetivos} específicos del proyecto.
		\item Identificar los riesgos que requieren una \textbf{mayor atención}.
	\end{itemize}
\end{itemize}
\begin{columns}
	\begin{column}{0.6\textwidth}
	\end{column}
	\begin{column}{0.4\textwidth}
		\includegraphics[width=3cm]{img/quantity}
	\end{column}
\end{columns}
}

\frame
{
\frametitle{Análisis Cuantitativo de Riesgos}
\begin{block}{Definición}
\begin{itemize}
	\item Analiza el efecto de esos riesgos y les asigna una \textbf{calificación numérica}.
	\item También presenta un método cuantitativo para \textbf{tomar decisiones} en caso de \emph{incertidumbre}. 
\end{itemize}
\end{block}
}


\frame
{
\frametitle{Análisis Cuantitativo de Riesgos}
\framesubtitle{Entradas}
\begin{columns}
	\begin{column}{0.8\textwidth}
		\begin{itemize}
			\item<1-> Registro de Riesgos.
			\item<2-> Plan de administración de \textbf{Riesgos}.
			\item<3-> Plan de administración de \textbf{Costos}.
			\item<4-> Plan de administración de \textbf{Planificación}.
			\item<5-> Activos organizacionales del Proceso.
		\end{itemize}
	\end{column}
	\begin{column}{0.2\textwidth}
		\includegraphics[width=2cm]{img/input}
	\end{column}
\end{columns}
}

\frame
{
\frametitle{Análisis Cuantitativo de Riesgos}
\framesubtitle{Herramientas y Técnicas}
\begin{columns}
	\begin{column}{0.8\textwidth}
		\begin{itemize}
			\item<1-> \emph{Recolección} de datos y técnicas de representación.
			\item<2-> Análisis cuantitativo de riesgos y técnicas de modelamiento.
			\begin{itemize}
				\item<3-> Análisis de sensibilidad.
				\item<4-> Análisis de valores de \emph{expectación monetaria}.
				\item<5-> Modelamiento  y simulación.
			\end{itemize}
			\item<6-> Juicio Experto.
		\end{itemize}
	\end{column}
	\begin{column}{0.2\textwidth}
		\includegraphics[width=2cm]{img/tools}
	\end{column}
\end{columns}
}

\frame
{
\frametitle{Análisis Cuantitativo de Riesgos}
\framesubtitle{Salidas}
\begin{columns}
	\begin{column}{0.8\textwidth}
		Actualización de los registros de riesgos.
		\begin{itemize}
			\item<1-> Análisis \textbf{probabilístico} del proyecto.
			\item<2-> Probabilidad de \emph{obtener} los objetivos de costo y tiempo.
			\item<3-> Lista \textbf{priorizada} de riesgos cuantificados.
			\item<4-> \textbf{Tendencias} en los resultados cuantitativos de análisis de riesgos.
		\end{itemize}
	\end{column}
	\begin{column}{0.2\textwidth}
		\includegraphics[width=2cm]{img/output}
	\end{column}
\end{columns}
}


\section{Respuestas al plan de riesgos}
\frame
{
\frametitle{Respuestas al plan de riesgos}

\begin{block}{Definición}
	Es el proceso de \textbf{desarrollo} de opciones y \textbf{oportunidades} para mejorar
	y reducir las \textbf{amenazas} a los objetivos del proyecto.
\end{block}
}

\frame
{
\frametitle{Respuestas al plan de riesgos}
\framesubtitle{Entradas}
\begin{columns}
	\begin{column}{0.8\textwidth}
		\begin{itemize}
			\item <1-> \textbf{Registro} de Riesgos
			\item <2-> \textbf{Plan} de Administración de Riesgos
		\end{itemize}
	\end{column}
	\begin{column}{0.2\textwidth}
		\includegraphics[width=2cm]{img/input}
	\end{column}
\end{columns}
}

%\frame
%{
%\frametitle{Respuestas al plan de riesgos}
%\framesubtitle{Herramientas y Técnicas}
%%Es necesario esto? quizas solo con el itemize basta.
%Varias estrategias de respuesta al riesgo están disponibles.
%La estrategia o combinación de estrategias más probable que sea eficaz
%ha de ser seleccionadas para cada riesgo. herramientas de análisis de
%riesgos, como el análisis de árbol de decisión, puede ser utilizado
%para elegir las respuestas más apropiadas.
%}

\frame
{
\frametitle{Respuestas al plan de riesgos}
\framesubtitle{Herramientas y Técnicas}

\begin{columns}
	\begin{column}{0.8\textwidth}
		\begin{itemize}
			\item <1-> Estrategias para Riesgos Negativos o \textbf{Amenazas}
				\begin{itemize}
					\item <2-> \textbf{Evitar}.
					\item <3-> \textbf{Transferencia}.
					\item <4-> \textbf{Mitigar}.
					\item <5-> \textbf{Aceptar}.
				\end{itemize}
			\item <6-> Estrategias para Riesgos Positivos u \textbf{Oportunidades}
				\begin{itemize}
					\item <7-> \textbf{Explotar}.
					\item <8-> \textbf{Compartir}.
					\item <9-> \textbf{Mejorar}.
					\item <10-> \textbf{Aceptar}.
				\end{itemize}
		\end{itemize}
	\end{column}
	\begin{column}{0.2\textwidth}
		\includegraphics[width=2cm]{img/tools}
	\end{column}
\end{columns}
}

\frame
{
\frametitle{Respuestas al plan de riesgos}
\framesubtitle{Herramientas y Técnicas}
\begin{columns}
	\begin{column}{0.8\textwidth}
		\begin{itemize}
			\item <1-> Estrategias de Respuestas de \textbf{Contingencias}
			\item <2-> Juicio \textbf{Experto}
		\end{itemize}
	\end{column}
	\begin{column}{0.2\textwidth}
		\includegraphics[width=2cm]{img/tools}
	\end{column}
\end{columns}
}

\frame
{
\frametitle{Respuestas al plan de riesgos}
\framesubtitle{Salidas}
\begin{columns}
	\begin{column}{0.8\textwidth}
		\begin{itemize}
			\item <1-> \textbf{Actualizaciones} de Registro de Riesgos
			\item <2-> Decisiones de \textbf{contrato} relacionadas a Riesgos
			\item <3-> \textbf{Actualizaciones} a Plan de Administración de Proyecto
			\item <4-> \textbf{Actualizaciones} al Documento de Proyecto
		\end{itemize}
	\end{column}
	\begin{column}{0.2\textwidth}
		\includegraphics[width=2cm]{img/output}
	\end{column}
\end{columns}
}

\section{Monitoreo y Control de Riesgos}
Monitorear y controlar los riesgos es el proceso de ejecución de planes de respuesta ante el riesgo, el seguimiento de los riesgos identificados, supervisar los riesgos residuales, identificando nuevos riesgos, y evaluar la eficacia del proceso de riesgo en todo el proyecto.

Respuestas de Planificación de riesgo que se incluyen en el plan de gestión de proyectos se ejecutan durante el ciclo de vida del proyecto, pero el trabajo del proyecto debe ser objeto de control permanente para las nuevos, cambiantes y actualizados riesgos.

El Monitor de Riesgos y Control de Procesos aplica técnicas, tales como la varianza y análisis de tendencias, que requieren el uso de información sobre el rendimiento generado durante la ejecución del proyecto:

\subsection{Entradas}

\begin{itemize}
	\item Registros de Riesgo
	\item Plan de Administración de Proyectos
	\item Información de Rendimiento en el Trabajo
	\item Reportes de Rendimiento
\end{itemize}

\subsection{Herramientas y Técnicas}

\begin{itemize}
	\item Revaloración de Riesgos
	\item Auditorías de Riesgos
	\item Varianza y Análisis de Tendencias
	\item Medición de Rendimiento Técnico
	\item Análisis de Reserva
	\item Reuniones de Estado
\end{itemize}

\subsection{Salidas}

\begin{itemize}
	\item Actualizaciones de Registro de Riesgos
	\item Actualizaciones de los Activos de Proceso de Organización
	\item Solicitudes de Cambio
	\item Actualizaciones al Plan de Administración de Proyecto
	\item Actualizaciones al Documento de Proyecto
\end{itemize}



\section{Conclusiones}
%APORTE RODRIGO:
La administración de los riesgos de un proyecto es un proceso que involucra a
todas las partes que o conforman, y que es muy necesario para poder estimar
los tiempos de desarrollo del mismo y de poder ayudar a evitar retrasos o el
fracaso del mismo. Además, las herramientas para poder realizar esto requieren de
personas que las manejen de buena forma y que puedan ver la visión holística
del problema para encontrar acciones de mitigación lo más adecuadas posibles.

%APORTE CRISTIAN:
Como equipo de trabajo tenemos muy clara, la importancia del todos los tópicos
que fueron repartidos entre los grupos de trabajos, pero creemos que la gestión
de los riesgos de un proyecto es fundamental, pues si no se tiene un especial cuidado
puedes llegar a acabar con un proyecto determinado.
En nuestra experiencia personal en la "Feria de Software", tuvimos la suerte de ser
los tres parte del mismo equipo, y al finalizar nos dimos cuenta que la parte más
dificil fué poder controlar los distintos riesgos de nuestro proyecto, tanto
del nivel técnico como a niver personal.

%APORTE GABRIEL:
Finalmente, se hace indispensable definir las respuestas que se pueden obtener al 
plantear el plan de riesgos, aplicando tanto herramientas, metodologías y 
técnicas, con el fin de manejar un conjunto de opciones a la hora de mitigar dichos
riesgos. Por otro lado todo proceso debe ser monitoreado y controlado, para ver 
que efectos tiene sobre el todo, con tal de hacer un seguimiento en el tiempo y
verificar que siga el camino establecido en los planes de gestión.

\end{document}
